\section{The ExtraTrees--TabPFN Model with Feature Augmentation and Heterogeneous Blending}\label{sec:v2-methods}

\subsection{Design Motivation}

To lift the accuracy of conditional-solubility prediction on unseen parents, this chapter walks through a progressive refinement spanning feature representation, learning criterion, and heterogeneous ensembling. Three objectives drive it: build high-dimensional explicit solute--solvent interaction features along the like-dissolves-like principle to break the ceiling of naive concatenation; inject the parent-equal prior into the tree-splitting criterion to correct long-tailed sampling bias; and bring in the pretrained tabular foundation model TabPFN v2, whose attention-based metric mechanism complements the orthogonal splits of decision trees through an out-of-fold non-negative convex blend.

\subsection{Base ExtraTrees Model}

Starting from the input defined in Section~\ref{sec:v2-common-problem}, the base model takes 11 physicochemical properties for the drug and 11 for the solvent, plus the measurement temperature, forming a 23-dimensional vector $[D_i; S_i; T_i]$.

For the learner we adopt Extremely Randomized Trees (ExtraTrees\cite{extratrees2006}). Unlike a standard random forest, ExtraTrees draws not only a random feature subset at each node but also fully random split thresholds within the feature's range, then keeps the split with the largest weighted variance reduction. This threshold randomisation pays off in cheminformatics: it smooths the regression staircase instead of fitting hard step boundaries, and it decorrelates the base learners, giving better variance control than a random forest under the collinearity of high-dimensional molecular descriptors. The base model uses 200 trees with row-equal weighting and serves as the reference for every later modification.

\subsection{Descriptors and Explicit Interaction Features}

To supply richer chemical information, the feature space grows from 11 base properties to the full 217 descriptors extracted by RDKit, spanning topological indices (BalabanJ, Chi), electronic structure (PEOE, partial charges), polar surface (LabuteASA, TPSA), and functional-group counts. Let the standardised drug and solvent descriptors be $D_i \in \mathbb{R}^{217}$ and $S_i \in \mathbb{R}^{217}$.

The next step addresses a fundamental shortcoming of the naive concatenation $[D_i; S_i; T_i]$ (435 dimensions): it fails to capture intermolecular interactions. In solution thermodynamics, dissolution reflects the competition among solute--solute, solvent--solvent, and solute--solvent interactions. Along the like-dissolves-like principle, the absolute-difference vector $|D_i - S_i|$ measures the mismatch between solute and solvent on matching physicochemical dimensions (polarity mismatch $|\mathrm{TPSA}_D - \mathrm{TPSA}_S|$, lipophilicity mismatch $|\mathrm{MolLogP}_D - \mathrm{MolLogP}_S|$; smaller means more compatible), while the Hadamard product $D_i \odot S_i$ measures the conjugated enhancement on specific thermodynamic properties (hydrogen-bond donor/acceptor interplay).

Leaving these interactions for the tree to discover implicitly would cost two or more orthogonal splits just to approximate a simple difference. We therefore build an 869-dimensional explicit interaction space:
\begin{equation}
x_i^{\mathrm{int}} = \left[ D_i;\; S_i;\; T_i;\; D_i \odot S_i;\; |D_i - S_i| \right] \in \mathbb{R}^{869}.
\end{equation}
Adding the standardised parent-scaffold descriptors yields a candidate 1086-dimensional representation. In the feature pipeline, median imputation and zero-variance filtering for the few malformed molecules are fitted strictly inside each inner training fold, leaving 665 dense, informative features and ruling out leakage.

\subsection{Parent-Weighted Optimisation}

To make ExtraTrees robust to long-tailed sampling, the parent weights $w_i(q)$ of Eq.~\eqref{eq:v2-train-weight} are plugged directly into the node-splitting algorithm. At a node $L$ the weighted sum of squared deviations is
\begin{equation}
J_L(v) = \sum_{i \in L} w_i(q) (y_i - v)^2.
\end{equation}
Setting $\frac{\partial J_L}{\partial v} = 0$ shows that the optimal leaf prediction is the weighted mean
\begin{equation}
\widehat v_L = \frac{\sum_{i \in L} w_i(q) y_i}{\sum_{i \in L} w_i(q)}.
\end{equation}
When a candidate split divides a parent $P$ into a left child $L$ and a right child $R$, its weighted impurity drop is
\begin{equation}
\Delta J(s) = J_P(\widehat v_P) - \left[ \frac{\sum_{i \in L} w_i}{\sum_{i \in P} w_i} J_L(\widehat v_L) + \frac{\sum_{i \in R} w_i}{\sum_{i \in P} w_i} J_R(\widehat v_R) \right].
\end{equation}
With $q=1$, every parent contributes the same total effective sample mass during tree growth: dense records inside a frequent parent can only refine local micro-splits and cannot dominate the rare parents at the macro level. The refined ExtraTrees uses 300 trees, a feature-sampling ratio of 0.7, a minimum leaf size of 2, and unrestricted depth.

\subsection{TabPFN Foundation Model}

As a heterogeneous complement to the tree ensemble, we incorporate the recent tabular foundation model TabPFN v2\cite{tabpfn2025}.

Pretrained via an autoregressive Transformer on millions of synthetic causal-topology prior datasets, TabPFN keeps its network parameters fixed at inference and takes the whole labelled training set as context:
\begin{equation}
\mathcal{C}_{\mathcal{T}} = \left\{ (x_i, y_i) \mid i \in \mathcal{T} \right\}.
\end{equation}
For a query condition $x^*$, it attends over all context samples and returns the posterior mean of the Bayesian predictive distribution:
\begin{equation}
f_{\mathrm{PFN}}(x^*) = \mathbb{E}_{\theta}\left[ Y \mid x^*, \mathcal{C}_{\mathcal{T}} \right].
\end{equation}
TabPFN's smooth attention-distance inductive bias captures gradual transitions in continuous space, whereas ExtraTrees builds staircase boundaries on high-dimensional sparse features through orthogonal splits. The two are highly independent in mechanism and in error distribution.

\subsection{Out-of-Fold Non-negative Convex Blending}

To blend ExtraTrees and TabPFN optimally without the second-order overfitting that comes from tuning the weight on the training set, we design an out-of-fold (OOF) non-negative convex combination. Inside the 3-fold split of the outer training set, let $t_i$ and $p_i$ be the ExtraTrees and TabPFN predictions of record $i$ on its held-out inner fold. The blended prediction is
\begin{equation}
\widehat y_i(\alpha) = (1 - \alpha) t_i + \alpha p_i = t_i + \alpha (p_i - t_i),
\end{equation}
where $(p_i - t_i)$ is TabPFN's correction to the tree baseline and $\alpha \in [0, 1]$ the blend weight. Substituting this into the parent-weighted risk (Eq.~\eqref{eq:v2-common-risk}) gives a constrained minimisation:
\begin{equation}
\widehat\alpha = \arg\min_{\alpha \in \{0, 0.25, 0.5, 0.75, 1.0\}} \frac{1}{|\mathcal{G}_{\mathrm{in}}|} \sum_{g \in \mathcal{G}_{\mathrm{in}}} \left( \frac{1}{|I_g|} \sum_{i \in I_g} \left| y_i - \left[ (1 - \alpha) t_i + \alpha p_i \right] \right| \right).
\label{eq:v2-blend-selection}
\end{equation}
A grid search picks the optimal step $\widehat\alpha$; $\widehat\alpha = 0$ collapses to the pure tree and $\widehat\alpha = 1$ to pure TabPFN. After the weight is chosen per outer fold, both learners are refit on the full outer data, and the final ensemble predicts
\begin{equation}
\widehat y(x) = (1 - \widehat\alpha) f_{\mathrm{ET}}(x) + \widehat\alpha f_{\mathrm{PFN}}(x).
\label{eq:v2-prediction-blend}
\end{equation}

\subsection{Ablation}

Holding the inner three-fold grouping and seed 2026 fixed on the development set, Table~\ref{tab:v2_tree_pilot} reports the progressive ablation of the core ExtraTrees designs.

\begin{table}[htbp]
\centering
\small
\caption{Ablation of the core ExtraTrees designs (values are parent-weighted MAE).}
\label{tab:v2_tree_pilot}
\begin{tabular}{lrr}
\toprule
Input representation & Row-equal ($q=0$) & Parent-equal ($q=1$) \\
\midrule
435-dim descriptor concatenation & 0.7112 & 0.6853 \\
869-dim explicit interactions (concat, product, absolute difference) & \textbf{0.6985} & \textbf{0.6760} \\
\bottomrule
\end{tabular}
\end{table}

The ablation confirms the incremental gain of each module. Relative to the plain concatenation, adding the absolute differences and products lowers the parent-weighted MAE by about 0.01, evidence that explicit interactions spare the tree from searching for them. Switching to parent-equal training brings a further reduction of over 0.022, confirming that the objective correction is necessary. On the TabPFN side, compressing the context to 8 or 16 records per parent worsens the MAE to above 0.74, while the full context lets the model exploit its large-sample background. The final out-of-fold two-branch blend again beats either single model, validating the complementarity of the heterogeneous ensemble.

Two practical issues were handled along the way. Non-finite descriptor values caused by dense heteroatoms are addressed with median imputation and dynamic variance filtering, fitted strictly inside each fold; the storage cost of TabPFN's large context is mitigated with dynamic memory release and GPU acceleration, bringing single inference below one second.

\subsection{Summary}

This chapter derived the ExtraTrees--TabPFN architecture and its progressive refinements. The experiments confirm that the 869-dimensional explicit physicochemical interactions, the parent-equal tree splitting, and the out-of-fold convex blending together build a heterogeneous ensemble that is both accurate and mechanistically complementary.
